\documentclass{article}
\usepackage{amsmath}
\usepackage[utf8]{inputenc}
\usepackage{booktabs}
\usepackage{microtype}
\usepackage[colorinlistoftodos]{todonotes}
\pagestyle{empty}

\title{Closest Pair Report}
\author{Author 1 and Author 2}

\begin{document}
  \maketitle

  \section{Results}

  The script did say that all my solutions were correct. At the largest input it takes:
  real	0m0.932s, user	0m0.950s, sys  0m0.027s
  The majority of that time likely comes from the recursive O(nlogn) logic and input parsing. The largest input has 1 000 000 inputs,
  so it makes somewhat sense.

  \section{Implementation details}

  \todo[inline]{How did you implement the solution? Which data structures were used? Which modifications to these data structures were used? What is the overall running time? Why?}

  I implemented the solution using a recursive algorithm. It was a divide and concour algorithm which
  used brute force for small inputs. The only data structures I used was a normal array for the initial list,
  and then a List<Vector2Int> inside of the logic. This could be made more efficiant for sure but .NET likely does
  a lot of optimization for me already.
  The overall running time is about 2-2.5 seconds for all inputs.

\end{document}
